\section{数値法}
\label{sec:numerics}

\subsection{何を刻むか}

内在角は閉形式で評価する。ルンゲ--クッタで \(\phi,\theta\) を進め、閉形式と照合するテストは、実装の退行を見るためだけに置く。本番経路は閉形式である。

時刻方向に刻んでよいのは、チャート上の力和を位置へ渡す作業段だけである。段幅 \(\Delta t\) は入力である。\(\ell\) と \(m\) からクーラン数を導出しない。

\subsection{拒否}

次のいずれかが起きた段は、状態を更新せず拒否を返す。

\begin{itemize}
\item いずれかのペアで \(r=0\)。
\item いずれかのペアで、使用した \(\gamma_s^{\mathrm{eq}}\) または \(\gamma_G,\gamma_E\) が零（許容 \(\varepsilon_\gamma\) 以内）。
\item 更新後の \texttt{TorsionParams} が \texttt{IsAdmissibleContinuous} を満たさない。
\item 契約ログが無い。
\end{itemize}

拒否のあとで許容錐へ射影しない。最後に成功した状態を、呼び出し側が保持する。ライブラリは自動で段幅を半分にして再試行しない。再試行は呼び出し側のループが、新しい契約ログとともに行う。

\subsection{保存の検査}

検査してよい保存量は、一粒子・一軸のヤコビ積分 \(E\) と \(p_\sigma\) である。閉形式をサンプルした値が、初期データの定数から許容誤差を超えて離れたら、そのサンプルを失敗にする。

多体チャートの運動エネルギーと、層力の仕事の和が保存する、という主張は DST の定理としてログに書かない。作業段がエネルギー差分を監視するなら、監視値の名前を \texttt{chart\_kinetic\_plus\_work} とし、\(E\) と加算しない。

\subsection{特異点の近く}

\(\gamma_s^{\mathrm{eq}}\) が許容の外側で小さいとき、力は大きくなる。その大きさを頭打ちにする平滑化は、記録の式を別の式に取り替える操作である。平滑化を入れるなら \texttt{working} の別関数にし、\texttt{gravityLayerTerm} のテスト期待値を平滑化後の数にしない。
